\section{Norms, maximum predecessors and exact counts}
\label{sec:descent}

All ordinals lie below $\varepsilon_0$ and are represented by hereditary
Cantor normal forms. Addition and multiplication by finite coefficients
have their ordinary ordinal meanings. For strictly decreasing exponents
and positive finite coefficients define
\[
 \norm(0)=0,\qquad
 \norm\!\left(\sum_{i<r}\omega^{\alpha_i}c_i\right)
       =\sum_{i<r}c_i(1+\norm(\alpha_i)).
\]
Thus $\norm(n)=n$, $\norm(\omega^d)=d+1$ for finite $d$,
$\norm(\omega^\omega)=3$, and
$\norm(\omega^{\omega+1})=4$. This is the additive norm used
in the relevant norm-controlled hierarchy framework
\cite{bcw,ww}.

\begin{lemma}[Effective norm balls]\label{lem:norm-balls}
Each norm ball is finite and effectively enumerable.
Cantor comparison and the norm are computable.
\end{lemma}
\begin{proof}
For budget zero there is only zero. In a nonzero ordinal of norm at
most $b$, every exponent has norm at most $b-1$. Inductively enumerate
those exponents, order them by recursive Cantor comparison, and enumerate
the finite coefficient lists whose total charge is at most $b$.
This produces every candidate and terminates. The same recursion
computes the norm and compares the resulting forms.
\end{proof}

Fix for now a computable strictly increasing $f:\N\to\N$ with
$f(n)>n$. Define
\[
 \psi_f(\alpha)=\max\bigl(\{0\}\cup
   \{\psi_f(\beta)+1:\beta<\alpha,\ 
                    \norm(\beta)\le f(\norm(\alpha))\}\bigr).
 \tag{2.1}
\]
Ordinal recursion is legitimate because the maximum is finite.
Its defining controlled implication is
\[
 \beta<\alpha,\quad \norm(\beta)\le f(\norm(\alpha))
 \quad\Longrightarrow\quad
 \psi_f(\beta)+1\le\psi_f(\alpha).
 \tag{2.2}
\]
Induction gives $\psi_f(n)=n$ for finite $n$.
No unrestricted monotonicity of $\psi_f$ is assumed.

\subsection{The exact maximum predecessor}
For $\alpha>0$ and $b\in\N$ put
\[
 M(\alpha,b)=\max\{\beta<\alpha:\norm(\beta)\le b\},
 \qquad L(e,b)=M(\omega^e,b).
\]
The candidate set contains zero. Put $M(0,b)=0$ as a convention.

\begin{theorem}[Greedy predecessor]\label{thm:greedy}
The following simultaneous recursion computes $M,L$ exactly.
Write $\alpha=\omega^e c+r$ with $c>0$, $r<\omega^e$, and
$w=1+\norm(e)$.
\begin{enumerate}[label=(\roman*)]
\item $M(\alpha,0)=0$ and $L(0,b)=L(e,0)=0$.
\item If $r>0$ and $cw\le b$, then
\[
 M(\alpha,b)=\omega^e c+M(r,b-cw).
\]
\item Otherwise, put
\[
 a=\begin{cases}
 \min(c-1,\lfloor b/w\rfloor),&r=0,\\
 \min(c,\lfloor b/w\rfloor),&r>0.
 \end{cases}
\]
Then $M(\alpha,b)=\omega^e a+L(e,b-aw)$.
\item If $e,b>0$, put $\eta=M(e,b-1)$,
$v=1+\norm(\eta)$ and $(q,s)=\operatorname{divmod}(b,v)$.
Then $L(e,b)=\omega^\eta q+L(\eta,s)$.
\end{enumerate}
Zero blocks are omitted.
\end{theorem}
\begin{proof}
When the whole leading block and a nonzero tail bound are affordable,
retaining that block and lowering the tail dominates every predecessor
with a smaller leading block. This proves (ii). Otherwise (iii) chooses
the largest affordable coefficient compatible with being below
$\alpha$; its remainder is strictly below $\omega^e$.

A nonzero ordinal below $\omega^e$ with budget $b$ has a leading
exponent $\eta'<e$ of norm at most $b-1$. Its largest possible
exponent is therefore $M(e,b-1)$. Cantor order next maximizes its
coefficient and then its lower tail, proving (iv).

The lexicographic measure (budget, phase), with $M$ above $L$,
proves termination. Clause (ii) lowers the budget. Clause (iii)
either lowers it or changes phase from $M$ to $L$. In (iv) the
first call uses $b-1$; also $v\le b$, so $q\ge1$ and $s<b$
for the second call. The value returned in each terminating case
is the greatest candidate just characterized.
\end{proof}

\begin{lemma}[Saturation]\label{lem:saturation}
For every nonzero limit $\lambda$ and every budget $b$,
$\norm(M(\lambda,b))=b$. The maximum is strictly increasing
with $b$.
\end{lemma}
\begin{proof}
If $\eta=M(\lambda,b)$ left unused budget
$a=b-\norm(\eta)>0$, the larger ordinal $\eta+a$ would still
be below the limit $\lambda$ and would have norm $b$, a
contradiction. If $b<c$, the ordinal $M(\lambda,b)+1$
is below $\lambda$ and has norm at most $c$, proving strict
increase.
\end{proof}

\begin{theorem}[Single-predecessor recurrence]\label{thm:predecessor}
For every nonzero limit $\lambda$ and $n\in\N$,
\[
 \psi_f(\lambda+n)
 =\psi_f\bigl(M(\lambda,f(\norm(\lambda)+n))\bigr)+1.
\]
\end{theorem}
\begin{proof}
Induct on $n$. Put $b=f(\norm(\lambda)+n)$ and
$\eta=M(\lambda,b)$. Saturation shows that $\eta$ is an admitted
predecessor of norm $b$, so it gives the lower inequality.
Every admitted $\beta<\lambda$ satisfies $\beta\le\eta$.
If $\beta<\eta$, then
$\norm(\beta)\le b=\norm(\eta)<f(\norm(\eta))$,
and controlled strictness bounds its contribution by
$\psi_f(\eta)$.

The other possible predecessors are $\lambda+j$ with $j<n$.
Put $b_j=f(\norm(\lambda)+j)$ and
$\eta_j=M(\lambda,b_j)$. Strict increase of $f$ and saturation
give $\eta_j<\eta$ and $\norm(\eta_j)=b_j<b$.
The induction hypothesis and controlled strictness yield
$\psi_f(\lambda+j)=\psi_f(\eta_j)+1\le\psi_f(\eta)$.
Thus each contribution to the defining maximum is at most
$\psi_f(\eta)+1$, proving the reverse inequality.
\end{proof}

\subsection{The complete event count}
For a nonfinite ordinal $\alpha$, write
$\alpha=\headpart(\alpha)+n$, where the first summand is a
nonzero limit and $n$ is finite. At each event replace
\[
 \alpha\longmapsto M(\headpart(\alpha),f(\norm(\alpha))).
 \tag{2.3}
\]
In particular the predecessor bound is the limit part, not the
original ordinal with its finite tail.

For a fixed limit $\lambda$ and total input norm
$t\ge\norm(\lambda)$, start at $\lambda+t-\norm(\lambda)$.
Let $\tau_\lambda(t)$ count events until the first finite
ordinal $Z_\lambda(t)$, and put
$J_\lambda(t)=\psi_f(\lambda+t-\norm(\lambda))$.

\begin{proposition}[Complete-count identity]\label{prop:complete}
The descent terminates and
\[
 Z_\lambda(t)=f^{\,\tau_\lambda(t)}(t),\qquad
 J_\lambda(t)=Z_\lambda(t)+\tau_\lambda(t).
\]
\end{proposition}
\begin{proof}
Every event strictly decreases the ordinal. Well-foundedness below
$\varepsilon_0$ therefore proves termination. Saturation gives the
successive norms $t,f(t),f^2(t),\ldots$; the finite endpoint equals
its own norm. Theorem~\ref{thm:predecessor} adds one for each
event, and finite evaluation returns $Z_\lambda(t)$.
\end{proof}

There is no further event count through finite successors below the
endpoint. A count stopped at an earlier nonfinite boundary is a
different segment count. The block calculations below always specify
which prefix remains and when their suffix first becomes finite.
